\chapter{de Sitter Horizons and Cosmological Probability}

\lectureoverview{A horizon is not merely a causal boundary. For a stationary observer it is a thermal system: a thermometer held a proper distance \(\rho\) from the horizon reads a temperature of order \(1/(2\pi\rho)\), while gravitational redshift makes the radiation seen far away much colder. We shall carry this black-hole lesson into de Sitter space, derive the fixed horizon radius \(H^{-1}\), and write the observer's static-patch metric. The second half of the lecture follows the thermodynamics to its unsettling consequences: finite de Sitter entropy, Poincare recurrences, Boltzmann brains, eternal inflation, and the measure problem. A final discussion separates the exact supersymmetric structures of string theory from the still-unknown framework that might describe our world.}

\section{What It Means for a Horizon to Be Hot}

We begin before the formal lecture, with the question that determines the evening's direction.

\begin{classroomqa}
\audiencequestion{A cosmic horizon contains no black hole. By what mechanism can it nevertheless be hot?\lecturetimestamp{00:00:17}}
\lecturerresponse{The mechanism is the same as for a black-hole horizon. The useful question is operational: what does a thermometer held near the horizon record, and what radiation reaches an observer far from it?}
\end{classroomqa}

The first experiment is local. Hold a thermometer at fixed proper distance \(\rho\) from any smooth, nonextremal horizon. Close enough to the horizon, the geometry has the universal Rindler form, and the thermometer reads
\begin{equation}
T_{\rm local}(\rho)\simeq \frac{1}{2\pi\rho}
\label{eq:l09-local-temperature}
\end{equation}
in units with \(\hbar=c=k_B=1\). Restoring constants gives
\begin{equation}
k_B T_{\rm local}(\rho)\simeq \frac{\hbar c}{2\pi\rho}.
\end{equation}
The divergence as \(\rho\to0\) does not mean that an infalling observer encounters an ordinary material furnace. It is the temperature assigned by an observer who remains stationary just outside the horizon and must therefore accelerate ever more violently.

The opening conversation also separates three notions that are easily conflated: energy, temperature, and entropy.

\begin{classroomqa}
\audiencequestion{If photons are sent into a black hole, do we still measure the same entropy, or does only the energy change?\lecturetimestamp{00:01:22}}
\lecturerresponse{Once a photon becomes inaccessible behind the horizon, additional information has been hidden. In the simplified bit language used here, one photon can hide one additional bit and thereby increase the entropy. Energy, temperature, and entropy are related, but they are not interchangeable.}
\end{classroomqa}

\editorialnote{For a physical black hole the entropy change is fixed thermodynamically by \(dS=dE/T_H\); it is not universally one bit for every photon, regardless of energy. The one-bit statement is the lecture's information-theoretic idealization.}

The near-horizon region continually emits and reabsorbs high-frequency quantum fluctuations. A horizon makes those fluctuations thermal for the stationary observer. Yet this raises an immediate puzzle. If the region is arbitrarily hot, why is a distant observer not blinded by arbitrarily energetic photons?

\section{Redshift and the Hawking Temperature}

An ordinary hot wall remains a source of high-frequency radiation even when we stand far from it. A horizon is different because its radiation must climb through a gravitational redshift. The photon's speed remains \(c\); what decreases is its energy,
\begin{equation}
E=\hbar\omega,
\qquad
\omega_{\rm received}<\omega_{\rm emitted}.
\end{equation}
The piece-of-chalk analogy used in the lecture is useful only up to this point. A thrown stone slows as it rises, whereas a photon continues at light speed while its frequency and energy decrease.

For a Schwarzschild black hole, the temperature measured by an observer at infinity is
\begin{equation}
T_H=\frac{1}{8\pi M}
\label{eq:l09-hawking-natural}
\end{equation}
in units \(G_N=\hbar=c=k_B=1\). In ordinary units,
\begin{equation}
T_H=\frac{\hbar c^3}{8\pi G_N M k_B}.
\label{eq:l09-hawking-full}
\end{equation}

\begin{figure}[htbp]
\centering
\includegraphics[width=0.68\textwidth]{lecture_09_figure_02.png}
\caption{The Hawking temperature \(T_H=1/(8\pi M)\) in natural units.}
\lectureframe{00:08:12}
\end{figure}

A larger black hole is colder. The lecture explains this by tracing a typical photon backward. Thermal photons have characteristic energy \(E_\gamma\sim T_H\). A photon that reaches infinity with energy of order \(M^{-1}\) had a much larger local energy when it began near the horizon. The redshift between the near-horizon layer and infinity is more severe for a larger hole, leaving a lower asymptotic temperature.

\begin{classroomqa}
\audiencequestion{If Schwarzschild time freezes at the horizon, should we follow the photon that originally fell in and built up the hole?\lecturetimestamp{00:04:19}}
\lecturerresponse{For the temperature question, follow the photons emitted outward. They begin at very high local frequency, lose energy while climbing out, and arrive as Hawking radiation. The fate of the original infalling photon is a separate information question.}
\end{classroomqa}

\section{The Planck Layer and the Stretched Horizon}

How close can the local relation in Equation~\eqref{eq:l09-local-temperature} be trusted? The lecture takes the Planck length as the shortest meaningful distance in the semiclassical description. This suggests a thin layer, roughly one Planck length outside the mathematical horizon, filled with Planck-scale fluctuations. It emits and absorbs continually.

Most emitted quanta do not escape. Only those directed into a very narrow outward cone avoid falling back. Quanta launched at generic angles return to the horizon. Outgoing and ingoing excitations collide, producing pairs and other particles; in this sense the stationary description contains a thin shell of extremely hot material.

\begin{figure}[htbp]
\centering
\includegraphics[width=0.74\textwidth]{lecture_09_figure_03.png}
\caption{The narrow escape cone for near-horizon radiation. Most locally emitted quanta return to the horizon.}
\lectureframe{00:11:38}
\end{figure}

There is a complementary spacetime picture. Away from a horizon, a virtual fluctuation may be represented schematically by a short closed loop: a particle-antiparticle excitation appears and disappears before it can become an asymptotic state. Near a horizon, a fluctuation can straddle the causal boundary. To the exterior observer, the outside portion looks like a particle emitted by the horizon and later reabsorbed. A very small fraction escapes and becomes Hawking radiation.

\begin{classroomqa}
\audiencequestion{In the exterior diagram, does such a fluctuation extend from \(t=-\infty\) to \(t=+\infty\)?\lecturetimestamp{00:15:00}}
\lecturerresponse{The exact classical horizon leads to that singular-looking description. In practice we replace it by a stretched horizon, a timelike surface about one Planck length outside. Emission and reabsorption then occur at finite exterior times.}
\end{classroomqa}

\begin{figure}[htbp]
\centering
\includegraphics[width=0.72\textwidth]{lecture_09_figure_04.png}
\caption{The classical horizon and the nearby stretched horizon used in the exterior description.}
\lectureframe{00:15:52}
\end{figure}

The stretched horizon is not a new invariant boundary. It is an effective surface placed where semiclassical geometry ceases to be reliable. Quantum fluctuations of matter and of the geometry itself smear the exact location on Planck scales. The construction packages that uncertainty into a surface on which the exterior description remains finite.

Nothing in this local argument knows whether the horizon belongs to a black hole or to an accelerating universe. That is the bridge to de Sitter space.

\section{de Sitter Space in Expanding Coordinates}

The formal lecture begins with the spatially flat slicing of de Sitter spacetime,
\begin{equation}
ds^2=-dt^2+e^{2Ht}d\mathbf{x}^2,
\qquad
d\mathbf{x}^2=dx^2+dy^2+dz^2.
\label{eq:l09-flat-slicing}
\end{equation}
The constant \(H\) is the de Sitter expansion rate. The factor appears as \(e^{2Ht}\) because the metric measures squared distance; a fixed coordinate interval has physical length
\begin{equation}
D(t)=e^{Ht}\lvert\Delta\mathbf{x}\rvert.
\end{equation}

\begin{figure}[htbp]
\centering
\includegraphics[width=0.78\textwidth]{lecture_09_figure_05.png}
\caption{The exponentially expanding de Sitter metric in flat coordinates.}
\lectureframe{00:18:50}
\end{figure}

These coordinates make expansion transparent but place future infinity infinitely far up the page. To bring it onto the blackboard, introduce conformal time
\begin{equation}
\eta=-\frac{1}{H}e^{-Ht},
\qquad
-\infty<\eta<0.
\label{eq:l09-conformal-time}
\end{equation}
Then
\begin{equation}
dt=\frac{d\eta}{-H\eta},
\qquad
e^{Ht}=-\frac{1}{H\eta},
\end{equation}
and Equation~\eqref{eq:l09-flat-slicing} becomes
\begin{equation}
ds^2=\frac{1}{H^2\eta^2}
\left(-d\eta^2+d\mathbf{x}^2\right).
\label{eq:l09-conformal-metric}
\end{equation}
As \(t\to+\infty\), \(\eta\to0^{-}\). Infinite proper time has been compressed to the finite horizontal line \(\eta=0\).

\editorialnote{The time transformation is corrected several times while being written on the board. Equation~\eqref{eq:l09-conformal-time} is the consistent form: early cosmic time maps to \(\eta\to-\infty\), and future infinity maps to \(\eta\to0^{-}\).}

\begin{figure}[htbp]
\centering
\includegraphics[width=0.82\textwidth]{lecture_09_figure_06.png}
\caption{The conformal-time transformation and the conformally flat metric.}
\lectureframe{00:21:24}
\end{figure}

The conformal factor multiplies both the temporal and spatial terms, so it does not alter null directions. For a radial light ray,
\begin{equation}
ds^2=0
\quad\Longrightarrow\quad
-d\eta^2+dx^2=0
\quad\Longrightarrow\quad
\frac{dx}{d\eta}=\pm1.
\label{eq:l09-null-rays}
\end{equation}
Light therefore travels at \(45^\circ\) in the conformal diagram. This is why the change of coordinates is useful: causal accessibility becomes visible.

\section{One Observer's Causal Patch}

Imagine an observer whose worldline continues for infinite proper time. On the compactified diagram it ends at one point on \(\eta=0\). Every signal the observer can ever receive lies inside the backward light cone of that endpoint. Events outside the cone can never communicate with the observer, however long the observer waits.

\begin{figure}[htbp]
\centering
\includegraphics[width=0.78\textwidth]{lecture_09_figure_07.png}
\caption{An observer's limiting backward light cone in conformal de Sitter coordinates.}
\lectureframe{00:24:38}
\end{figure}

A second observer has a different cone. The two accessible regions overlap, but neither is the whole spacetime. For the remainder of the argument we adopt the first observer's point of view and call that backward light cone the observer's causal patch.

The drawing seems to make the patch shrink toward the future, but coordinate width is not proper distance. On a constant-\(\eta\) slice the squared physical distance between the observer at \(x=0\) and the horizon at \(x=\Delta x\) is
\begin{equation}
D^2=\frac{(\Delta x)^2}{H^2\eta^2}.
\label{eq:l09-horizon-distance-start}
\end{equation}
The horizon is a \(45^\circ\) null line ending at \(\eta=0\), hence
\begin{equation}
\Delta x=\lvert\eta\rvert.
\end{equation}
Substitution gives
\begin{equation}
D_{\rm hor}^2=\frac{1}{H^2},
\qquad
D_{\rm hor}=\frac{1}{H}.
\label{eq:l09-horizon-radius}
\end{equation}
The coordinate separation shrinks, but the conformal factor grows by exactly the compensating amount. The proper horizon radius is constant.

\begin{figure}[htbp]
\centering
\includegraphics[width=0.80\textwidth]{lecture_09_figure_08.png}
\caption{The board derivation that converts the shrinking coordinate interval into the fixed proper distance \(H^{-1}\).}
\lectureframe{00:28:42}
\end{figure}

\section{The Static Patch}

A fixed horizon radius suggests coordinates adapted to the central observer. Without deriving the coordinate transformation, the lecture writes the result:
\begin{equation}
ds^2=
-\left(1-H^2r^2\right)d\tau^2
+\frac{dr^2}{1-H^2r^2}
+r^2d\Omega_2^2.
\label{eq:l09-static-patch}
\end{equation}
The coefficients depend on \(r\), not on \(\tau\). This chart covers the observer's static patch,
\begin{equation}
0\le r<\frac{1}{H}.
\end{equation}

\begin{figure}[htbp]
\centering
\includegraphics[width=0.88\textwidth]{lecture_09_figure_09.png}
\caption{The de Sitter static-patch metric and its comparison with the Schwarzschild factor.}
\lectureframe{00:34:22}
\end{figure}

The resemblance to Schwarzschild is immediate:
\begin{equation}
f_{\rm dS}(r)=1-H^2r^2,
\qquad
f_{\rm Schw}(r)=1-\frac{2G_NM}{r}.
\end{equation}
In either case the static-coordinate horizon occurs where \(f(r)=0\). For de Sitter,
\begin{equation}
1-H^2r^2=0
\quad\Longrightarrow\quad
r_H=H^{-1}.
\end{equation}
The chart becomes singular there just as Schwarzschild coordinates become singular at a black-hole horizon. The global geometries are nevertheless opposite in an important sense: the black-hole observer stands outside and looks inward, whereas the de Sitter observer stands inside and looks outward.

\begin{classroomqa}
\audiencequestion{How can the geometry be static when freely released particles move apart in an exponentially expanding universe?\lecturetimestamp{00:36:06}}
\lecturerresponse{Static does not mean that every particle remains fixed. It means that the metric coefficients and the outcome of a repeated experiment are independent of the start time. Release two particles now and they accelerate apart; repeat the same experiment later and the same motion occurs.}
\end{classroomqa}

\begin{classroomqa}
\audiencequestion{Do the two separating particles see one another redshift?\lecturetimestamp{00:38:34}}
\lecturerresponse{Yes. Their increasing separation produces a relative redshift, even though the observer-adapted metric is time independent.}
\end{classroomqa}

\section{Temperature Seen from Inside}

The near-horizon structures of Schwarzschild and de Sitter space are locally equivalent. The de Sitter horizon therefore emits and reabsorbs thermal quanta. A rare photon propagates inward to the central observer, losing energy as it climbs the static-patch potential. The centrally measured Gibbons--Hawking temperature is
\begin{equation}
T_{\rm dS}=\frac{H}{2\pi}
\label{eq:l09-de-sitter-temperature}
\end{equation}
in natural units. A stationary thermometer at radius \(r\) reads the Tolman-redshifted value
\begin{equation}
T_{\rm local}(r)=
\frac{T_{\rm dS}}{\sqrt{1-H^2r^2}},
\label{eq:l09-tolman}
\end{equation}
which diverges as \(r\to H^{-1}\) and reduces near the horizon to Equation~\eqref{eq:l09-local-temperature}.

\editorialnote{The lecture obtains the scale from the horizon distance and redshift rather than writing Equations~\eqref{eq:l09-de-sitter-temperature} and \eqref{eq:l09-tolman} explicitly. They state the standard normalization of that argument.}

For the observed late-time acceleration, the lecture uses the period estimate
\begin{equation}
H^{-1}\sim 20\times10^9\ {\rm light\ years}.
\label{eq:l09-lecture-horizon-estimate}
\end{equation}
A typical horizon photon then has a wavelength comparable to billions of light years. This is not microwave-background radiation; it is vastly colder.

The operational thought experiment is a fishing pole on a cosmological scale. Alice remains at \(r=0\), attaches a thermometer to a very long line that also transmits data, and lets cosmic repulsion carry the instrument outward. She stops it just short of the horizon, waits for the signal to return, and finds a large local temperature. The thermometer beside her at the center reads the tiny value \(H/(2\pi)\). One horizon supports both readings because they are made at different gravitational potentials.

\section{Crossing the Horizon}

The classical outside description resembles the familiar black-hole story. Bob at the center watches Alice recede. Her signals become increasingly redshifted and arrive increasingly slowly; she appears flattened against the horizon and ultimately fades. A later traveler, Jane, appears as another thinner layer in a sedimentary stack near the boundary.

\begin{classroomqa}
\audiencequestion{From our point of view, do objects ever pass through the cosmic horizon?\lecturetimestamp{00:45:30}}
\lecturerresponse{Not in the static observer's classical description. Their signals redshift and fade while their apparent position approaches the horizon. In their own freely falling description, however, they cross without encountering a locally singular surface.}
\end{classroomqa}

\begin{classroomqa}
\audiencequestion{Is the horizon hot only for observers who remain close without crossing? An infalling observer is not cooked merely by passing through it, is she?\lecturetimestamp{00:48:27}}
\lecturerresponse{Alice crosses smoothly in her own frame. The extreme temperature belongs to the stationary description. If Bob insists on resolving Alice ever closer to the horizon, he must illuminate her with ever shorter wavelengths. The attempted Heisenberg microscope then deposits enough energy to cook what it tries to observe.}
\end{classroomqa}

This is not a contradiction between two local events. The stationary observer and the infaller use radically different measurements. Horizon complementarity asks how both descriptions can be valid without allowing one observer to compare duplicate information directly.

\section{A Newtonian Model of Vacuum Repulsion}

At this point a student asks what replaces the black hole's gravitational well. Susskind answers by constructing a Newtonian model. Let \(\phi\) be the gravitational potential, so the acceleration of a test particle is
\begin{equation}
\mathbf{a}=-\boldsymbol{\nabla}\phi.
\end{equation}
For ordinary nonrelativistic matter, Poisson's equation is
\begin{equation}
\nabla^2\phi=4\pi G_N\rho.
\label{eq:l09-poisson}
\end{equation}

\begin{classroomqa}
\audiencequestion{If the black-hole effects come from its gravitational well, what is the corresponding source in de Sitter space?\lecturetimestamp{00:49:36}}
\lecturerresponse{The vacuum energy itself supplies a spatially uniform gravitational source. In a Newtonian description it produces a quadratic potential rather than a \(1/r\) well.}
\end{classroomqa}

\begin{figure}[htbp]
\centering
\includegraphics[width=0.78\textwidth]{lecture_09_figure_10.png}
\caption{Poisson's equation and the lecture's schematic radial second-derivative form.}
\lectureframe{00:51:38}
\end{figure}

The board suppresses normalization, spherical-measure terms, pressure, and initially even the overall sign, reducing the exercise to
\begin{equation}
\frac{d^2\phi}{dr^2}\sim H^2.
\end{equation}
The essential inference is that a constant source gives a quadratic function. With the sign required for de Sitter repulsion,
\begin{equation}
\phi_{\rm dS}(r)=-\frac{1}{2}H^2r^2,
\qquad
a_r=-\frac{d\phi_{\rm dS}}{dr}=H^2r.
\label{eq:l09-de-sitter-potential}
\end{equation}
The acceleration points outward and grows linearly with distance. Solving \(\ddot r=H^2r\) gives
\begin{equation}
r(\tau)=A e^{H\tau}+B e^{-H\tau},
\end{equation}
so the growing branch reproduces exponential separation.

\begin{figure}[htbp]
\centering
\includegraphics[width=0.76\textwidth]{lecture_09_figure_11.png}
\caption{The uniform-source equation and the quadratic de Sitter potential.}
\lectureframe{00:53:32}
\end{figure}

\editorialnote{In general relativity the weak-field acceleration of a homogeneous fluid depends on \(\rho+3p\), not on \(\rho\) alone. Vacuum energy has \(p=-\rho\), so \(\rho+3p=-2\rho\) and the acceleration is repulsive. Equation~\eqref{eq:l09-de-sitter-potential} is the correctly normalized de Sitter result; the board calculation intentionally tracks only the quadratic dependence and sign.}

Place an ordinary mass \(M\) at the center. The qualitative potential becomes
\begin{equation}
\phi(r)=-\frac{G_NM}{r}-\frac{1}{2}H^2r^2,
\label{eq:l09-combined-potential}
\end{equation}
and the radial acceleration is
\begin{equation}
a_r=-\frac{G_NM}{r^2}+H^2r.
\end{equation}
At small \(r\), local attraction wins; at large \(r\), vacuum repulsion wins. Their balance occurs at
\begin{equation}
r_\ast=\left(\frac{G_NM}{H^2}\right)^{1/3}.
\label{eq:l09-turnaround}
\end{equation}
This is the clean version of the lecture's potential maximum. It explains qualitatively why nearby gravitationally bound systems can remain together while sufficiently distant galaxies accelerate away.

\begin{figure}[htbp]
\centering
\includegraphics[width=0.78\textwidth]{lecture_09_figure_12.png}
\caption{Ordinary \(1/r\) attraction combined with the upside-down quadratic de Sitter potential.}
\lectureframe{00:58:20}
\end{figure}

\begin{classroomqa}
\audiencequestion{Do ordinary forces still matter, or does vacuum repulsion eventually tear apart every island of matter?\lecturetimestamp{00:57:29}}
\lecturerresponse{All forces must be included. On sufficiently small scales, electromagnetic binding and local gravity dominate. Vacuum repulsion becomes important only across enormous distances set by \(H^{-1}\) and by the balance scale appropriate to the bound mass.}
\end{classroomqa}

\section{One Source, Several Names}

The discussion now pauses over terminology. Vacuum energy, a cosmological constant, and dark energy are not three additive substances in the model being used. They are three descriptions of the same \(w=-1\) source. One may speak of a constant term in Einstein's equation, a Lorentz-invariant energy density of the vacuum, or the component driving late-time acceleration.

\begin{classroomqa}
\audiencequestion{Are the uniform Newtonian source and the quantum vacuum energy both present as separate effects?\lecturetimestamp{00:59:14}}
\lecturerresponse{No. The uniform density is the Newtonian model of the same source. Vacuum energy, cosmological constant, and, in this lecture, dark energy refer to one physical component.}
\end{classroomqa}

\begin{classroomqa}
\audiencequestion{Does the same phenomenon therefore have a classical description and a quantum description?\lecturetimestamp{01:00:32}}
\lecturerresponse{Its microscopic origin may be quantum mechanical, just as part of an electron's mass has a quantum origin. Once the effective energy density is known, it enters the classical gravitational equations as an ordinary source. We do not know the ultimate origin of the observed value.}
\end{classroomqa}

In Planck units, the observed vacuum-energy density is extraordinarily small,
\begin{equation}
\rho_\Lambda\sim10^{-123}\rho_{\rm Pl},
\label{eq:l09-planck-smallness}
\end{equation}
up to the convention used for the Planck density. That small number is the cosmological-constant problem. It is also why the horizon is so large: in four dimensions,
\begin{equation}
H^2=\frac{8\pi G_N}{3}\rho_\Lambda,
\qquad
r_H=H^{-1}.
\end{equation}
If \(H\) were larger, the horizon would be smaller, the central de Sitter temperature would be higher, and the repulsive acceleration \(H^2r\) would be stronger.

\begin{classroomqa}
\audiencequestion{Is the tiny cosmological constant related to Planck's constant and quantum fluctuations?\lecturetimestamp{01:03:27}}
\lecturerresponse{The clean comparison uses dimensionless Planck units, where \(G_N=\hbar=c=1\). In those units the vacuum density is about \(123\) orders of magnitude below the natural Planck density. Explaining that hierarchy is one of the central puzzles of fundamental physics.}
\end{classroomqa}

\begin{classroomqa}
\audiencequestion{Must vacuum energy be an exactly constant cosmological constant? Could the accelerated component decrease and eventually settle to zero?\lecturetimestamp{01:04:52}}
\lecturerresponse{That is a logically possible alternative: a dynamical dark-energy field could vary with time. The lecture adopts the constant case because observations available at the time increasingly favored behavior close to a cosmological constant, not because every alternative had been proved impossible.}
\end{classroomqa}

\section{Astronomy in the Distant Future}

In an asymptotically de Sitter universe, structures that are not gravitationally bound eventually cross the observer's horizon. The cosmic microwave background redshifts below detectability. Distant galaxies disappear. A surviving astronomer may see only a merged local gravitational island inside an otherwise dark patch.

\begin{classroomqa}
\audiencequestion{Without a telescope, would anyone even notice that the distant galaxies had disappeared?\lecturetimestamp{01:06:17}}
\lecturerresponse{Probably not. The serious problem is deeper: even with advanced instruments, future observers may lack the distant tracers and microwave background from which we infer expansion today.}
\end{classroomqa}

The old island-universe picture could then seem compelling. Olbers' paradox would not force a different conclusion because it assumes an infinite, static, uniform population of luminous sources; the future observer would see only the finite stars of the local island. Dark energy does not continually manufacture replacement galaxies from empty space.

The horizon in exact de Sitter space does not shrink around the local galaxy. Its physical radius remains \(H^{-1}\), enormously larger than a galaxy. What disappears is the unbound matter carried beyond that fixed causal boundary.

\begin{classroomqa}
\audiencequestion{Could stars ejected from the galaxy serve as test bodies for detecting the large-scale geometry?\lecturetimestamp{01:11:25}}
\lecturerresponse{For a while, yes. Ejected objects and long-baseline light signals could probe the geometry. But the effect becomes appreciable only over distances of order a substantial fraction of a billion light years, so the experiment would take cosmological time.}
\end{classroomqa}

\begin{classroomqa}
\audiencequestion{How do we know a distant horizon has not already hidden most of the universe from us?\lecturetimestamp{01:11:45}}
\lecturerresponse{We already see enough galaxies, redshifts, and background radiation to measure the expansion and reconstruct the large-scale geometry. Future observers deprived of those tracers would face a much harder inference problem.}
\end{classroomqa}

\begin{classroomqa}
\audiencequestion{If the horizon lay just beyond Andromeda, would its presence be obvious?\lecturetimestamp{01:12:19}}
\lecturerresponse{Yes. Such a small horizon would require a much larger \(H\), and the resulting repulsion would strongly alter Andromeda's motion. The measured value instead places the asymptotic horizon on a vastly larger scale.}
\end{classroomqa}

\begin{classroomqa}
\audiencequestion{Could new big bangs be occurring beyond our horizon?\lecturetimestamp{01:13:25}}
\lecturerresponse{Events beyond the horizon are not visible to us or to our descendants. A related but more precise possibility, bubble nucleation in eternal inflation, will be introduced later in the lecture.}
\end{classroomqa}

With rockets, clocks, and several billion years, future observers could send explorers across a large fraction of the patch and reconstruct its metric. On ordinary laboratory or civilizational time scales, however, the effect is almost inaccessible. This distinction corrects the overly strong claim that the expansion would be impossible to discover: it is operationally remote, not forbidden in principle.

\begin{classroomqa}
\audiencequestion{Could future explorers determine the geometry much sooner than waiting for a thermal recurrence?\lecturetimestamp{01:21:40}}
\lecturerresponse{Certainly. A recurrence is not the proposed measurement. Explorers exchanging light signals across billion-light-year baselines could map the static patch on a vastly shorter, though still cosmological, time scale.}
\end{classroomqa}

\begin{classroomqa}
\audiencequestion{Would redshift erase all information returned by such explorers?\lecturetimestamp{01:23:24}}
\lecturerresponse{Not if they probe scales well inside the horizon. A baseline of order a billion light years is already sensitive to expansion while producing only a moderate redshift. The experiment is limited chiefly by its immense duration.}
\end{classroomqa}

\section{Finite Entropy and Recurrence}

The static patch is not empty in the thermodynamic sense. Its horizon has area
\begin{equation}
A_H=4\pi H^{-2},
\end{equation}
temperature \(T_{\rm dS}=H/(2\pi)\), and entropy
\begin{equation}
S_{\rm dS}=\frac{A_H}{4G_N}
=\frac{\pi}{G_NH^2}
\label{eq:l09-de-sitter-entropy}
\end{equation}
in units \(\hbar=c=k_B=1\). The entropy is finite. One may therefore picture the static patch, with suitable caution, as a finite-entropy thermal cavity whose excitations are concentrated near the horizon in the stationary description.

An ordinary box of gas at equilibrium appears featureless on short time scales, but statistical mechanics does not forbid fluctuations. If a macrostate has an entropy deficit \(\Delta S\) relative to equilibrium, its probability is roughly
\begin{equation}
P\sim e^{-\Delta S}.
\label{eq:l09-fluctuation-probability}
\end{equation}
The corresponding waiting time grows as
\begin{equation}
t\sim t_0 e^{\Delta S}.
\end{equation}
For a finite system, the time to revisit a state arbitrarily close to an earlier one is the Poincare recurrence time,
\begin{equation}
t_{\rm rec}\sim H^{-1}e^{S_{\rm dS}}.
\label{eq:l09-recurrence}
\end{equation}
Since \(S_{\rm dS}\) for our horizon is of order \(10^{120}\), the lecture summarizes the scale as \(e^{10^{120}}\). Whether the prefactor is measured in years, Hubble times, or microscopic units is irrelevant beside the double exponential.

\begin{figure}[htbp]
\centering
\includegraphics[width=0.76\textwidth]{lecture_09_figure_13.png}
\caption{The recurrence discussion, with the enormous exponential waiting-time scale on the board.}
\lectureframe{01:16:24}
\end{figure}

Before such a recurrence, smaller fluctuations are vastly more common. The molecules in a room might all gather in one corner; a complicated vortex might briefly form; in the lecture's playful hierarchy, thermal motion might assemble an elephant. The smaller the required entropy deficit, the cheaper the fluctuation.

\begin{classroomqa}
\audiencequestion{Are Poincare recurrences specifically a consequence of quantum mechanics?\lecturetimestamp{01:25:44}}
\lecturerresponse{No. Recurrence and rare fluctuations already arise in classical statistical mechanics. Their application to de Sitter space uses our present physical framework, which might someday be revised, but thermodynamics and statistical reasoning are among the most robust structures we know.}
\end{classroomqa}

The meaning of probability itself may deepen in a future theory. Susskind does not claim otherwise. His narrower point is that abandoning statistical mechanics merely because its long-time implications are strange would discard one of physics' best-tested organizing principles.

\section{Boltzmann Brains and Typical Observers}

Suppose a thermal fluctuation creates an entire big-bang-like state. That requires an enormous entropy decrease. A fluctuation producing one solar system is cheaper; an already warm Earth is cheaper still. If the only requirement is one observer, the least costly fluctuation might create only a functioning brain complete with apparent memories. Such a hypothetical observer is called a Boltzmann brain.

The issue is not whether one bizarre fluctuation is logically possible. It is whether a cosmological model predicts that almost every observer is such a fluctuation. We reason from our observations by assuming that we are not fantastically atypical. A theory in which ordinary observers with coherent cosmic histories are exponentially outnumbered by isolated thermal brains undermines that inference.

The Boltzmann-brain problem therefore names a quantitative consistency test. One compares the expected abundance of ordinary observers with the abundance of fluctuation-born observers. If the latter dominate, either the vacuum must decay before they proliferate, the measure must be changed, or the cosmological assumptions must be reconsidered.

An earlier question about an ultraviolet--infrared connection is explicitly deferred. The recurrence argument does not depend on resolving that separate issue.

\section{Eternal Inflation and Bubble Universes}

The lecture next introduces, without attempting a full derivation, a second probability problem. A metastable de Sitter vacuum can nucleate a bubble in which the vacuum energy is lower. The bubble wall expands, and the interior can develop into an open Friedmann--Robertson--Walker universe. Meanwhile the surrounding inflating region keeps expanding and nucleates more bubbles.

\begin{figure}[htbp]
\centering
\includegraphics[width=0.74\textwidth]{lecture_09_figure_14.png}
\caption{A lower-vacuum-energy bubble being drawn inside an inflating de Sitter background.}
\lectureframe{01:29:38}
\end{figure}

The crucial feature is exponential population growth. New physical volume is produced faster than bubbles remove it, so the number of bubble regions grows without bound. Different bubbles may settle into different low-energy states and therefore exhibit different effective physics.

To expose the statistical difficulty, Susskind first leaves cosmology and considers an exponentially growing human population with finite but unknown resources. Assume that observers are typical among all people who ever live. If the population doubles every generation and then ends, half of all people live in the last generation, three quarters in the last two, and seven eighths in the last three:
\begin{equation}
f_n=1-2^{-n}.
\label{eq:l09-last-generations}
\end{equation}
Typicality would then place us uncomfortably near the end, even though we know nothing about the total resource supply. This is not offered as a reliable prediction of human extinction. It is a warning that exponential growth makes counts dominated by the cutoff.

The coin-flipper example states the ordinary typicality assumption more cleanly. If many experimenters each flip a fair coin \(N\) times, the overwhelming majority obtain
\begin{equation}
n_H=\frac{N}{2}+O(\sqrt{N}).
\end{equation}
An extreme observer should first suspect a bad model or a bad coin rather than assume without evidence that she occupies a fantastically rare tail. Cosmological probability uses the same typicality instinct, but in an infinite population it needs a well-defined way to count.

\section{Escher's Boundary and the Measure Problem}

Susskind's geometric analogy is M.\ C.\ Escher's \emph{Circle Limit IV}, the disk tiled by angels and devils that appear smaller toward the boundary. In the hyperbolic geometry represented by the drawing, the figures are equivalent; their apparent Euclidean size changes only because of the projection.

\begin{figure}[htbp]
\centering
\includegraphics[width=0.76\textwidth]{lecture_09_figure_15.png}
\caption{The blackboard surrogate for Escher's exponentially proliferating figures near a circular boundary.}
\lectureframe{01:37:54}
\end{figure}

Imagine that Escher has finite time, ink, and patience. Near the edge he eventually stops, perhaps leaving some figures incomplete. Because the number of figures grows exponentially with hyperbolic radius, a finite outer layer contains a fixed-order fraction of every figure drawn. Roughly half can lie in the final layer, no matter how far outward the cutoff is moved.

Now ask for the ratio of angels to devils. In the ideal infinite pattern symmetry suggests one-to-one equality. In any finite count, however, a slightly different boundary can include many more figures of one type. Since most of the count lies near the boundary, a small deformation of that boundary produces an order-one change in the answer.

\begin{figure}[htbp]
\centering
\includegraphics[width=0.76\textwidth]{lecture_09_figure_16.png}
\caption{The cutoff boundary dominates the count in an exponentially growing pattern.}
\lectureframe{01:43:06}
\end{figure}

The same ambiguity appears in eternal inflation. Count blue and pink bubble universes below some time cutoff, then push the cutoff later. Both counts diverge, and their ratio can depend sensitively on what \emph{time} means and how the cutoff surface is drawn. A probability measure is a prescription that regulates the infinite population and extracts relative probabilities. The failure of different plausible prescriptions to agree is the measure problem.

\begin{classroomqa}
\audiencequestion{Is there a relation between this cutoff problem and renormalization?\lecturetimestamp{01:45:43}}
\lecturerresponse{There is an analogy. In quantum field theory one introduces a lattice or ultraviolet cutoff, refines it, and seeks quantities insensitive to the regulator's edge details. Cosmologists hope for an equally robust extraction from eternal inflation, but no comparably decisive prescription is yet available in the framework discussed here.}
\end{classroomqa}

The analogy should not be overstated. Renormalization succeeds when a controlled flow identifies regulator-independent observables. The measure problem is precisely that the analogous invariant questions and limiting procedure remain unclear. Saying that we may be asking the wrong question is not a solution until the right question is supplied.

\begin{classroomqa}
\audiencequestion{Is this bubble picture the same phenomenon as inflation in our early universe?\lecturetimestamp{01:48:19}}
\lecturerresponse{It is closely connected. In eternal inflation, the parent background continues inflating while local bubble interiors undergo their own cosmological evolution. Our observable universe could lie inside one such region.}
\end{classroomqa}

\section{Could Bubble Collisions Be Seen?}

Two bubbles do not simply add their vacuum energies when they meet. A domain wall separates regions with different vacuum states, and the geometry of the wall determines what either side can observe. Depending on the collision's location relative to our past light cone, its imprint may be hidden or visible.

\begin{classroomqa}
\audiencequestion{What happens when two vacuum bubbles collide?\lecturetimestamp{01:52:07}}
\lecturerresponse{The collision creates a domain-wall geometry whose behavior depends on the two vacua and on the wall tension. Under favorable causal circumstances, the disturbance can enter our observable region and leave an anisotropic signature.}
\end{classroomqa}

The lecture discusses a cold region in the cosmic microwave background as a possible candidate and points to polarization around a circular feature as a discriminant. This is deliberately presented as a possibility, not a detection. A convincing collision signature would be revolutionary because it would supply direct evidence for neighboring bubble regions and make the measure problem experimentally urgent.

\begin{figure}[htbp]
\centering
\includegraphics[width=0.70\textwidth]{lecture_09_figure_17.png}
\caption{A circular sky feature and its possible polarization pattern in the bubble-collision discussion.}
\lectureframe{01:50:22}
\end{figure}

\editorialnote{The cold-spot discussion records the observational status and proposals at the time of the lecture. A temperature anomaly by itself is not evidence of a bubble collision; a detailed, correlated signature would be required.}

\begin{classroomqa}
\audiencequestion{Are the bubbles four-dimensional, and what spatial curvature do their interiors have?\lecturetimestamp{01:53:37}}
\lecturerresponse{They are spacetime regions, not three-dimensional soap bubbles evolving in an external time. Their interiors are typically open FRW universes with negative spatial curvature. A collision perturbs that geometry but does not simply turn negative curvature into positive curvature.}
\end{classroomqa}

\section{What String Theory Does and Does Not Establish}

The final discussion begins with an invitation to assign a probability that string theory is correct. Susskind refuses the number because the proposition is too ambiguous. He instead separates two claims.

First, there is a highly constrained mathematical structure combining quantum mechanics and gravity. It has generated precise results and deep mathematical identities, and its internal consistency has survived extensive scrutiny. In the lecture's informal vocabulary this is \emph{String Theory} with a capital \(S\). The versions understood most sharply rely heavily on supersymmetry.

Second, the observed world is not supersymmetric at accessible energies. Connecting the exact supersymmetric constructions to our nonsupersymmetric universe requires supersymmetry breaking and a larger framework that is not controlled with the same precision. In the lecture's vocabulary, that hoped-for extension is \emph{string theory} with a small \(s\).

\begin{classroomqa}
\audiencequestion{Can one assign a probability that string theory is the correct theory of nature?\lecturetimestamp{01:54:18}}
\lecturerresponse{Not responsibly without specifying which claim is meant. The exact mathematical structures exist and include quantum gravity, but the rigorously controlled supersymmetric models are not literally our observed world. Whether a larger extension contains our world is a matter on which physicists can have informed prejudices, not a well-defined numerical probability.}
\end{classroomqa}

The analogy is with mechanics before the full mathematical machinery was available. One might know \(F=ma\) and the inverse-square force yet be able to solve only highly symmetric circular orbits. The solvable sector would be genuine but not the whole theory. Likewise, exact supersymmetric string constructions may be a special, unusually tractable region of a larger structure.

\begin{figure}[htbp]
\centering
\includegraphics[width=0.60\textwidth]{lecture_09_figure_18.png}
\caption{The lecture's nested picture: rigorously controlled supersymmetric string constructions inside a conjecturally larger framework.}
\lectureframe{01:57:34}
\end{figure}

\begin{classroomqa}
\audiencequestion{Why not simply break supersymmetry in the exact theory?\lecturetimestamp{01:57:02}}
\lecturerresponse{Many approximate constructions study supersymmetry breaking, but the fully controlled nonsupersymmetric backbone is precisely what is missing. One can calculate after making standard approximations without always knowing the exact object those approximations approach.}
\end{classroomqa}

\begin{classroomqa}
\audiencequestion{Does the exact supersymmetric theory rule out a larger, less symmetric extension that contains our world?\lecturetimestamp{02:00:33}}
\lecturerresponse{No. It establishes that the observed world is not described by the exact supersymmetric sector alone. It neither proves nor disproves that this sector is connected to a larger theory containing realistic physics.}
\end{classroomqa}

\begin{classroomqa}
\audiencequestion{Do string constructions include symmetry breaking analogous to particle-physics symmetry breaking?\lecturetimestamp{02:01:21}}
\lecturerresponse{There is extensive work on supersymmetry breaking, and much is known about supersymmetric mathematics. The unresolved point is a precise, nonperturbative formulation that combines realistic breaking with the full theory rather than only with controlled approximations.}
\end{classroomqa}

\begin{classroomqa}
\audiencequestion{Is special relativity postulated in string theory, or does it emerge as a prediction?\lecturetimestamp{02:03:02}}
\lecturerresponse{The underlying rules are Lorentz invariant, even though a particular background full of matter can select a preferred rest frame. Gravity is still more unavoidable: the string spectrum contains a massless spin-two excitation with the interactions of a graviton.}
\end{classroomqa}

\begin{classroomqa}
\audiencequestion{Is string theory consistent with general relativity and background independence?\lecturetimestamp{02:04:11}}
\lecturerresponse{Its low-energy gravitational dynamics are consistent with general relativity, and local backgrounds can be changed dynamically. The subtlety lies in asymptotic boundary conditions: they label the problem itself and cannot generally be changed by a local operation within that spacetime.}
\end{classroomqa}

The lecture closes on that distinction between a robust mathematical core and an unfinished physical extrapolation. It is not a verdict that the core is useless, nor a claim that the extrapolation is established. It is a statement of where exact knowledge ends.
